\section{Persistence}
\label{sec:Persistence}

The Archimesh server persists all model state to Neo4j. This chapter describes
the graph model, repository layer and key invariants.

\subsection{Repository layer}

The persistence layer is organized into a split repository with a main
coordinator and specialized support classes:

\begin{description}
\item[\texttt{Neo4jRepositoryImpl}] Main coordinating repository. Orchestrates
  op-log writes, materialized state mutations, head revision updates,
  snapshots, rebuilds, imports, exports, compactions and integrity checks.
\item[\texttt{Neo4jOpLogSupport}] Commit and operation persistence. Writes
  \texttt{:Commit} and \texttt{:Op} nodes, maintains the \texttt{NEXT}
  chain.
\item[\texttt{Neo4jMaterializedStateSupport}] Materialized graph mutation.
  Creates, updates and deletes elements, relationships, views, view objects,
  connections and folders. Applies LWW merges and tombstone checks.
  Maintains property clocks.
\item[\texttt{Neo4jReadSupport}] Snapshot export, checkout delta reads,
  consistency checks. Reads materialized state and converts it to the JSON
  snapshot/export format.
\item[\texttt{Neo4jCompactionSupport}] Metadata compaction. Prunes old
  commits and property clocks below the retention watermark. Reports
  eligibility of tombstones and field clocks for future compaction.
\end{description}

\subsection{Graph model}

The database stores four kinds of state per model:

\begin{enumerate}
\item Model catalog and ACL metadata
\item Materialized current state
\item Committed op-log history
\item Merge metadata (tombstones, property clocks)
\end{enumerate}

All identity is model-scoped: materialized entity identity is
\texttt{(modelId, id)}; commit identity is \texttt{(modelId, opBatchId)};
tag identity is \texttt{(modelId, tagName)}.

\subsubsection{Core model node}

Every model is rooted at a \texttt{:Model} node:

\begin{lstlisting}[caption={Model node properties}]
{
  modelId, modelName,
  registered: true,
  headRevision: <number>,
  createdAt, updatedAt,
  adminUsers: ["user1", ...],
  writerUsers: ["user2", ...],
  readerUsers: ["user3", ...]
}
\end{lstlisting}

ACL data is stored directly on \texttt{:Model} as user-id arrays.

\subsubsection{Materialized state nodes}

\begin{longtable}{p{30mm}p{124mm}}
\toprule
Label & Properties \\
\midrule
\texttt{:Element} & \texttt{modelId}, \texttt{id}, \texttt{archimateType},
  \texttt{name}, \texttt{documentation},
  \texttt{name\_lamport}, \texttt{name\_clientId},
  \texttt{documentation\_lamport}, \texttt{documentation\_clientId} \\
\texttt{:Relationship} & \texttt{modelId}, \texttt{id}, \texttt{archimateType},
  \texttt{name}, \texttt{documentation}, \texttt{sourceId}, \texttt{targetId},
  plus LWW clocks for name, documentation, source and target. \\
\texttt{:Folder} & \texttt{modelId}, \texttt{id}, \texttt{folderType},
  \texttt{name}, \texttt{documentation}, plus LWW clocks. \\
\texttt{:View} & \texttt{modelId}, \texttt{id}, \texttt{name},
  \texttt{documentation}, \texttt{notationJson}, plus LWW clocks. \\
\texttt{:ViewObject} & \texttt{modelId}, \texttt{id}, \texttt{notationJson},
  plus per-field LWW clocks for geometry, style and label fields. \\
\texttt{:Connection} & \texttt{modelId}, \texttt{id},
  \texttt{sourceViewObjectId}, \texttt{targetViewObjectId},
  \texttt{representsId}, \texttt{notationJson}, plus per-field LWW clocks. \\
\bottomrule
\end{longtable}

\subsubsection{Materialized relationships}

\begin{lstlisting}[caption={Core relationships}]
(:Model)-[:HAS_ELEMENT]->(:Element)
(:Model)-[:HAS_REL]->(:Relationship)
(:Relationship)-[:SOURCE]->(:Element)
(:Relationship)-[:TARGET]->(:Element)

(:Model)-[:HAS_FOLDER]->(:Folder)
(:Folder)-[:HAS_FOLDER]->(:Folder)
(:Folder)-[:CONTAINS_ELEMENT]->(:Element)
(:Folder)-[:CONTAINS_REL]->(:Relationship)
(:Folder)-[:CONTAINS_VIEW]->(:View)

(:Model)-[:HAS_VIEW]->(:View)
(:View)-[:CONTAINS]->(:ViewObject)
(:ViewObject)-[:REPRESENTS]->(:Element)
(:ViewObject)-[:CHILD_MEMBER]->(:ViewObject)

(:View)-[:CONTAINS]->(:Connection)
(:Connection)-[:REPRESENTS]->(:Relationship)
(:Connection)-[:FROM]->(:ViewObject)
(:Connection)-[:TO]->(:ViewObject)
\end{lstlisting}

\subsubsection{Folder rules}

\begin{itemize}
\item Root folders use stable synthetic ids:
  \texttt{folder:root-strategy}, \texttt{folder:root-business},
  \texttt{folder:root-diagrams}, etc.
\item User folders use immutable ids --- rename changes the name property
  but not the node identity.
\item A model-tree object belongs to exactly one folder at a time.
\item Cross-root folder moves are rejected by validation.
\item Non-empty folder delete is rejected.
\end{itemize}

\subsubsection{Op-log history}

Commit nodes and operation nodes form the append-only operation log:

\begin{lstlisting}[caption={Op-log nodes}]
(:Commit {
  modelId, opBatchId,
  revisionFrom, revisionTo,
  userId, sessionId, ts
})

(:Op {
  seq, type, targetId, payloadJson
})
\end{lstlisting}

Relationships:

\begin{lstlisting}[caption={Op-log relationships}]
(:Model)-[:HAS_COMMIT]->(:Commit)
(:Commit)-[:HAS_OP]->(:Op)
(:Commit)-[:NEXT]->(:Commit)
\end{lstlisting}

\begin{itemize}
\item \texttt{revisionFrom..revisionTo} is the assigned revision range for
  one committed submit batch.
\item \texttt{payloadJson} stores the normalized op payload.
\item \texttt{seq} preserves op order inside a batch.
\item \texttt{NEXT} keeps the timeline linear.
\end{itemize}

\subsubsection{Tags}

Tags are stored as standalone nodes:

\begin{lstlisting}[caption={Tag node}]
(:ModelTag {
  modelId, tagName, description,
  revision, createdAt, snapshotJson
})
\end{lstlisting}

\begin{itemize}
\item Tags are immutable --- once created, they cannot be modified or deleted
  (tag deletion is disabled by default).
\item \texttt{snapshotJson} stores the materialized snapshot at the tagged
  revision.
\item Tags do not create branches; they are read-only pointers on the linear
  timeline.
\end{itemize}

\subsubsection{Merge metadata --- tombstones}

Tombstones prevent stale recreates after deletes:

\begin{lstlisting}[caption={Tombstone nodes}]
(:ElementTombstone { modelId, id, lamport, clientId })
(:RelationshipTombstone { modelId, id, lamport, clientId })
(:ViewTombstone { modelId, id, lamport, clientId })
(:ViewObjectTombstone { modelId, id, lamport, clientId })
(:ConnectionTombstone { modelId, id, lamport, clientId })
\end{lstlisting}

Tombstones are anchored from \texttt{:Model} via specific relationship types
(\texttt{HAS\_ELEMENT\_TOMBSTONE}, etc.). Compaction reports tombstone
eligibility but retains them for delete safety.

\subsubsection{Merge metadata --- property clocks}

Property merge state is stored in dedicated metadata nodes:

\begin{lstlisting}[caption={PropertyClock node}]
(:PropertyClock {
  modelId, ownerId, key,
  lamport, clientId, deleted
})
\end{lstlisting}

Anchored via \texttt{(:Model)-[:HAS\_PROPERTY\_CLOCK]->(:PropertyClock)}.
Supports LWW and OR-set style property semantics. Compaction may prune old
property clocks below the retention watermark.

\subsection{Neo4j constraints}

The schema file at \filepath{neo4j/schema.cypher} defines the following
uniqueness constraints:

\begin{lstlisting}[caption={Uniqueness constraints}]
(:Model {modelId})                    UNIQUE
(:Element {modelId, id})              UNIQUE
(:Relationship {modelId, id})         UNIQUE
(:View {modelId, id})                 UNIQUE
(:Folder {modelId, id})               UNIQUE
(:ViewObject {modelId, id})           UNIQUE
(:Connection {modelId, id})           UNIQUE
(:Commit {modelId, opBatchId})        UNIQUE
(:ModelTag {modelId, tagName})        UNIQUE
\end{lstlisting}

An index on \texttt{(:Commit)} by \texttt{(modelId, revisionFrom, revisionTo)}
supports efficient delta window queries.

\subsection{Persistence invariants}

\begin{itemize}
\item Materialized entity identity is \texttt{(modelId, id)} --- never rely
  on node-internal ids across queries.
\item Commit/idempotency identity is \texttt{(modelId, opBatchId)} ---
  duplicate submission of the same batch is a no-op.
\item Repository write paths must fail fast on Neo4j persistence errors ---
  never silently skip a mutation.
\item Notation field definitions must flow through
  \texttt{NotationMetadata}, the shared source of truth for supported
  notation fields and persisted merge metadata.
\item Compaction never allows stale \texttt{Create*} or \texttt{Update*} to
  beat a logically newer delete --- tombstones are retained.
\end{itemize}

\subsection{Configuration}

\begin{longtable}{p{42mm}p{50mm}p{62mm}}
\toprule
Property & Env var & Default \\
\midrule
\texttt{app.neo4j.uri} & \texttt{NEO4J\_URI} &
  \texttt{bolt://localhost:7687} \\
\texttt{app.neo4j.username} & \texttt{NEO4J\_USER} & \texttt{neo4j} \\
\texttt{app.neo4j.password} & \texttt{NEO4J\_PASSWORD} &
  \texttt{devpassword} \\
\bottomrule
\end{longtable}

In production, always set \texttt{NEO4J\_PASSWORD} to a strong value. The
default \texttt{devpassword} is for local development only.
